% 第9章：同一联合分布的两种分解（教学构造，数值可逐格复算）。
% 两种邮件 a（普通通知）、b（含链接促销），y=1 表示诈骗。
% 三月（源域）：P^X(a)=0.9, P^X(b)=0.1；四月（目标域）：P^X(a)=0.1, P^X(b)=0.9。
% 两个月都有 eta(a)=P(y=1|a)=0.1, eta(b)=0.9。
% 联合概率 P(x,y)=P^X(x)eta：三月 (a,0)=.81 (a,1)=.09 (b,0)=.01 (b,1)=.09；
% 四月 (a,0)=.09 (a,1)=.01 (b,0)=.09 (b,1)=.81。
% 类别比例：三月 P^Y(1)=.18，四月 .82；类内分布 P(x|y)=P(x,y)/P^Y(y)：
% 三月 P(a|0)=.81/.82=81/82, P(b|0)=1/82, P(a|1)=P(b|1)=1/2；
% 四月 P(a|0)=1/2, P(a|1)=1/82, P(b|1)=81/82。
\begingroup
\definecolor{FigureGreen}{HTML}{829C8B}
\definecolor{FigureGreenFill}{HTML}{E8EEE9}
\newcommand{\ltNineCell}[5]{%
  \node[draw=FigureGrid,fill=#5,minimum width=#3 mm,minimum height=6.4mm,
    inner sep=0pt,anchor=north west,font=\figurefont\footnotesize]
    at (#1,#2) {#4};}
\begin{tikzpicture}[figure picture,x=1mm,y=1mm]
\foreach \row/\top/\name/\pa/\pb/\jaz/\jao/\jbz/\jbo/\pyz/\pyo/\az/\bz/\ao/\bo in {%
  0/60/三月（源域）/0.9/0.1/0.81/0.09/0.01/0.09/0.82/0.18/{81/82}/{1/82}/{1/2}/{1/2},%
  1/27/四月（目标域）/0.1/0.9/0.09/0.01/0.09/0.81/0.18/0.82/{1/2}/{1/2}/{1/82}/{81/82}}{%
  \node[figure heading,anchor=west] at (0,\top+4.5) {\name};
  % 左块：输入比例与条件标签概率
  \ltNineCell{0}{\top}{12}{$x$}{white}
  \ltNineCell{12}{\top}{15}{$\mathcal P^X(x)$}{white}
  \ltNineCell{27}{\top}{15}{$\eta(x)$}{white}
  \ltNineCell{0}{\top-6.4}{12}{$a$}{white}
  \ltNineCell{12}{\top-6.4}{15}{\pa}{FigureRedFill}
  \ltNineCell{27}{\top-6.4}{15}{0.1}{FigureGreenFill}
  \ltNineCell{0}{\top-12.8}{12}{$b$}{white}
  \ltNineCell{12}{\top-12.8}{15}{\pb}{FigureRedFill}
  \ltNineCell{27}{\top-12.8}{15}{0.9}{FigureGreenFill}
  % 中块：联合分布
  \ltNineCell{62}{\top}{12}{}{white}
  \ltNineCell{74}{\top}{14}{$y=0$}{white}
  \ltNineCell{88}{\top}{14}{$y=1$}{white}
  \ltNineCell{62}{\top-6.4}{12}{$a$}{white}
  \ltNineCell{74}{\top-6.4}{14}{\jaz}{FigureBlueFill}
  \ltNineCell{88}{\top-6.4}{14}{\jao}{FigureBlueFill}
  \ltNineCell{62}{\top-12.8}{12}{$b$}{white}
  \ltNineCell{74}{\top-12.8}{14}{\jbz}{FigureBlueFill}
  \ltNineCell{88}{\top-12.8}{14}{\jbo}{FigureBlueFill}
  % 右块：类别比例与类内输入分布
  \ltNineCell{122}{\top}{11}{$y$}{white}
  \ltNineCell{133}{\top}{14}{$\mathcal P^Y(y)$}{white}
  \ltNineCell{147}{\top}{17}{$\mathcal P(a\mid y)$}{white}
  \ltNineCell{122}{\top-6.4}{11}{0}{white}
  \ltNineCell{133}{\top-6.4}{14}{\pyz}{FigureRedFill}
  \ltNineCell{147}{\top-6.4}{17}{\az}{FigureRedFill}
  \ltNineCell{122}{\top-12.8}{11}{1}{white}
  \ltNineCell{133}{\top-12.8}{14}{\pyo}{FigureRedFill}
  \ltNineCell{147}{\top-12.8}{17}{\ao}{FigureRedFill}
  % 箭头：联合分布分别分解
  \draw[figure flow] (61,\top-9.6) -- (43,\top-9.6);
  \draw[figure flow] (103,\top-9.6) -- (121,\top-9.6);
}
\node[figure note] at (21,73.5) {输入比例 $\times$ 条件标签};
\node[figure note] at (82,73.5) {联合分布 $\mathcal P(x,y)$};
\node[figure note] at (143,73.5) {类别比例 $\times$ 类内分布};
% 图例
\node[draw=FigureGrid,fill=FigureGreenFill,minimum width=4mm,minimum height=3mm,inner sep=0pt] at (30,-2) {};
\node[figure note,anchor=west] at (33,-2) {两个月相同};
\node[draw=FigureGrid,fill=FigureRedFill,minimum width=4mm,minimum height=3mm,inner sep=0pt] at (70,-2) {};
\node[figure note,anchor=west] at (73,-2) {两个月不同};
\node[figure note,anchor=west] at (100,-2) {$\mathcal P(b\mid y)=1-\mathcal P(a\mid y)$};
\end{tikzpicture}
\endgroup
